% p3-appA-formal-defs

\section{P3 Appendix A: Formal Definitions}

Appendix A: Formal Definitions

            Figure (fig-p3-appA-formal). Formal-definitions triptych --- Pellis-formal / Trinity-formal /
                                                 pi-L-formal.

  A.1 Pellis Expansion: Formal Definition
  Definition A.1 (Pellis sequence). A Pellis sequence for x $\in$ R>0 is any sequence (ck)k $\geq$
  0 $\subseteq$ Z$\geq$ 0 such that x = sumk=0$\in$fty ck $\varphi$-k, converging in R. The standard Pellis
  sequence uses the floor function (Algorithm 4.1). The centred Pellis sequence uses the
  nearest-integer function, allowing signed coefficients and potentially faster
  convergence for certain targets.

  Lemma A.2 (Finiteness). For x $\in$ Q($\varphi$) = Q(5), the standard Pellis sequence is eventually
  zero (termination in finitely many steps). For x notin Q(5), it has infinitely many nonzero
  terms.

  A.2 Trinity Monomial: Formal Definition
  Definition A.3 (Trinity monomial). A Trinity monomial is any element of T = { n $\cdot$ 3k $\cdot$ $\varphi$p
  $\cdot$ pim $\cdot$ eq : n $\in$ Z$\geq$ 1, (k,p,m,q) $\in$ Z4 }.

  Definition A.4 (Complexity function). The complexity function ell : Z$\geq$ 1 $\times$ Z4 $\rightarrow$ Z$\geq$ 0 is
  as in Section 2.6.

  Definition A.5 (Trinity distance). The Trinity distance from x to T at budget L is dL(x) =
  minT $\in$ T, ell(lambdaT) $\leq$ L |x - T|.

  A.3 Projection Map: Formal Definition
  Definition A.6 (Projection map with tiebreaking). The projection map is PiL : R>0 $\rightarrow$ T,
  PiL(x) = argminT $\in$ T, ell(lambdaT) $\leq$ L |x - T|, with tiebreaking rule: in case of a tie,
  choose the monomial with smallest n, then smallest |k|+|p|+|m|+|q|, then
  lexicographically among labels of equal total absolute value.

  A.4 AlphaPhi Objects: Canonical Definitions
  Definition A.7 (Algebraic AlphaPhi). alpha$\varphi$( alg) = $\varphi$-3/2 = (5-2)/2 approx 0.11803. This
  is an element of Q(5); its Pellis expansion terminates.

  Definition A.8 (Logarithmic AlphaPhi). alpha$\varphi$(log) = 2ln$\varphi$/pi approx 0.30635. This is
  transcendental; its Pellis expansion does not terminate.

  Note: alpha$\varphi$( alg) $\neq$ alpha$\varphi$(log); the two objects play distinct roles and must not be
  conflated. The notation inconsistency in Chapter 38 of the PhD program is superseded
  by these definitions.

  A.5 Operator Definitions
  Definition A.9 (Pellis shift). On X = T $\times$ ZN:
  sigma(T, (c0, c1, c2, …)) = (PiL(T - c0), (c1, c2, c3, …)).

  Definition A.10 (Koopman operator). For f $\in$ L2(X, mu): K f = f circ sigma.

  Definition A.11 (Transfer operator). For weight w and f $\in$ C(X): (Lw f)(x) = sumy :
  sigma(y) = x w(y) f(y).

\subsection{References}

- Vasilev, D. \textit{Flos Aureus Monograph: Trinity Constants}. DOI \href{https://zenodo.org/doi/10.5281/zenodo.19227877}{10.5281/zenodo.19227877}.
- CODATA 2018 recommended values: NIST \href{https://physics.nist.gov/cuu/Constants/}{https://physics.nist.gov/cuu/Constants/}.
- Heyrovska, R. Golden-angle FSC publications --- Heyrovský Institute of Physical Chemistry.
